\chapter{The I/O Page}\label{cha:io}

One page, \texttt{\$D000-\$DFFF}, always visible. The authoritative
register-level reference is the machine's own headers, and
Appendix~\ref{cha:registers} is generated from them at build time --- the
way the Neo6502 handbook generates its keyword reference from the
firmware source. This chapter is the map of what lives where:

\begin{center}
\small
\begin{tabular}{@{}l l >{\raggedright\arraybackslash}p{0.60\linewidth}@{}}
\toprule
\texttt{\$D000} & VICKY   & video: layers, palette, blitter, sprites, SHEILA\\
\texttt{\$D100} & input   & keyboard FIFO and its status, the held keys, the mouse, break-pending\\
\texttt{\$D200} & DMA     & copy / fill / swap / line / triangle\\
\texttt{\$D300} & storage & the host filesystem: open, read, dir, chdir\ldots\\
\texttt{\$D480} & OPL2    & the YM3812: address, data, status --- the AdLib's arrangement\\
\texttt{\$D500} & system  & clock, RTC, frame counter, DUMP, the settings the ROM reads, the sound sequencer at \texttt{\$D5E0}\\
\texttt{\$D600} & banks   & the eight bank registers, masks at \texttt{\$D620}\\
\texttt{\$D700} & MATH    & IEEE floats, transcendentals, math lists, mul/div\\
\texttt{\$D800} & Tube    & the co-processor: status, byte in, byte out, start/stop\\
\texttt{\$D900} & N:      & the network: four channels, tcp://, http(s):// and tnfs:// (see \texttt{core/net.h})\\
\texttt{\$DA00} & JIM     & the terminal: a VT100/ANSI in hardware, drawing on the console's screen (see \texttt{core/term.h})\\
\texttt{\$DF00} & far gate& call slots, descriptor table pointer, depth\\
\bottomrule
\end{tabular}
\end{center}

\section{A split screen: SHEILA}\label{sec:split}
The Apple~II put four lines of text under its graphics; C64 programmers
did the same with a raster interrupt timed to the line. On this machine
the CPU does nothing at all: SHEILA, VICKY's display-list coprocessor
(the Amiga's copper is its ancestor), runs a short list from memory at
every frame, in step with the raster, and a register it writes takes
effect on the line about to be drawn. A split is six instructions:

\begin{type}
MOVE $10,$00     line 0: layer 0, the console's text, off
MOVE $20,$19     layer 1, a 640-wide 8 bpp bitmap, on
WAIT 416         26 text rows of 16 lines: the split
MOVE $20,$00     the bitmap off
MOVE $10,$07     the text on: rows 26-29 are text
END              and again from the top, next frame
\end{type}

\noindent Each instruction is four bytes --- the operation (\texttt{1}
WAIT, \texttt{2} MOVE, \texttt{0} END), then its operands --- and the
list is anywhere in memory, a 28-bit address in \texttt{\$D060}--\texttt{\$D063},
started by \texttt{1} in \texttt{\$D064}. A MOVE is exactly a write by
the CPU, so any VICKY register can change mid-frame: the mode, the
scroll, a layer, the palette. Two rules: line numbers count the glass,
0--479, whatever the mode (in a 240-line mode each line is drawn twice,
so split on an even one); and \textbf{nothing turns SHEILA off after a
program}, so a program that turns it on turns it off when it leaves.

\screen{shots/split.png}{\texttt{SPLIT}: blitter lines above, the
console's last four rows below, and SHEILA holding the line between
them.}

\verb|SPLIT| is the demonstration, and it has two twins that draw the
same ribbon of lines --- endpoints bouncing, the line sixteen behind
erased --- as fast as they can, and say how fast in the text band:
\pth{/LANG/EHBASIC/EX/SPLIT.BAS}, with EhBASIC's own \verb|LINE|, and
\pth{/LANG/MSBASIC/SPLIT.BAS}, where Microsoft BASIC, which has no
graphics words, sets the blitter's registers with \verb|POKE|. At the
40.5\,MHz a host starts at:

\begin{center}
\small
\begin{tabular}{@{}l r >{\raggedright\arraybackslash}p{0.48\linewidth}@{}}
\toprule
\textbf{Language} & \textbf{lines/s} & \textbf{how it draws one} \\
\midrule
C (\verb|SPLIT|) & 7400 & fifteen register writes \\
EhBASIC & 330 & \verb|LINE|, one blitter operation \\
Microsoft BASIC & 140 & ten \verb|POKE|s in a subroutine \\
\bottomrule
\end{tabular}
\end{center}

\noindent The blitter draws a line in no time at all, so the numbers
are the languages: what a BASIC spends is its own arithmetic, its
arrays and its line lookups, and three hundred lines a second is a
good deal more than a game of 1985 asked for. The split itself costs
every one of them nothing. A faster clock (F7, or \verb|SETUP|) scales
all three.
